\subsection{Payday Loans and Credit Cards}

\subsubsection{Credit Cards}

\content{
        \begin{example} Credit cards usually compound interest daily. Suppose
        that you owe \$1000 on a credit card with a 25\% APR. You make a payment
        of \$100 (after interest is compounded for the month). How much will you
        still owe the bank? (Assume each month has exactly 30 days)
        \end{example}
}{% presentation
\vspace{3in}
}{% nopresentation
\vspace{3in}
}{% comments
Use the compount interest formula, where the exponent is just 30. This might
take some explaining.
}

\content{
        \begin{example} Suppose that you owe \$1000 on a credit card with 28\%
        APR, and the minimum payment each month is \$50, additionally assume
        that you make no additional purchases (not a very realistic assumption,
        I know, but let's make this simple for ourselves).
        How much will we owe the bank after 3 months? 
        \end{example}
}{% presentation
\vspace{4in}
}{% nopresentation
\vspace{4in}
}{% comments
}

\content{
        \begin{example} Let's go the computer for this one. In the above
        situation from the previous example, how long will it take us to pay off
        the bank? How much do we pay? 
        \end{example}
}{% presentation
}{% nopresentation
}{% comments
}

\subsubsection{Payday Loans}

\content{
        \begin{example} Payday loans are short-term loans where the interest is
        calculated up front. Consider the following example. You need \$400
        dollars to make rent, but don't get paid for another two weeks. You
        decide to take out a payday loan, and the lender charges \$15 dollars
        interest for every \$100 dollars borrowed. How much interest will you
        pay if you pay back the loan on time? What is the APR of this loan?
        \end{example}
}{% presentation
\vspace{3in}
}{% nopresentation
\vspace{3in}
}{% comments
}


